\documentclass[UTF8,a4paper,11pt,fontset=windows]{ctexart}
\usepackage[margin=2.15cm,headheight=15pt]{geometry}
\usepackage{amsmath,amssymb,mathtools,booktabs,longtable,array,tabularx,enumitem,xcolor,fancyhdr,hyperref}
\definecolor{ink}{HTML}{18374D}
\definecolor{accent}{HTML}{167D8D}
\definecolor{soft}{HTML}{EAF2F5}
\hypersetup{colorlinks=true,linkcolor=ink,urlcolor=accent,pdftitle={CMC真题考点与教材阅读重点},pdfauthor={学习资料整理}}
\pagestyle{fancy}\fancyhf{}
\fancyhead[L]{\small CMC 数学类：真题考点与教材阅读重点}
\fancyhead[R]{\small 材料分析}
\fancyfoot[C]{\small\thepage}
\setlength{\parindent}{2em}
\setlength{\parskip}{4pt}
\setlist[itemize]{leftmargin=1.6em,itemsep=3pt,topsep=3pt}
\setlist[enumerate]{leftmargin=1.8em,itemsep=3pt,topsep=3pt}
\renewcommand{\arraystretch}{1.25}
\newcolumntype{Y}{>{\raggedright\arraybackslash}X}
\newcolumntype{P}[1]{>{\raggedright\arraybackslash}p{#1}}
\newcommand{\key}[1]{\textbf{\color{ink}#1}}
\newcommand{\para}[1]{\par\medskip\noindent\key{#1}\quad}
\newcommand{\topic}[1]{\section*{#1}\addcontentsline{toc}{section}{#1}}
\pdfstringdefDisableCommands{\def\quad{ }}
\newcommand{\exref}[1]{\par\noindent{\small\color{accent}真题定位：#1}\par}
\newcommand{\R}{\mathbb R}
\newcommand{\C}{\mathbb C}
\newcommand{\dd}{\,\mathrm d}
\newcommand{\rank}{\operatorname{rank}}
\newcommand{\tr}{\operatorname{tr}}
\newcommand{\Span}{\operatorname{span}}
\newcommand{\priority}[1]{\textbf{#1}}
\begin{document}
\begin{center}
{\Huge\bfseries\color{ink}CMC 真题考点}\\[7pt]
{\LARGE\bfseries\color{ink}与教材阅读重点}\\[8pt]
{\large 数学专业类，第 1--17 届真题汇编与现有重排教程对照}\\[5pt]
{\small 整理日期：2026 年 10 月 23 日}
\end{center}

\topic{先看这页：应该把时间花在哪里}
本文以你提供的 225 页解析版为依据，逐卷查看 60 套试卷的题目，建立 503 个题位的分类索引，再对照工作区内两册教程及数学分析各章重排源码。\key{503 是题位数，含 A/B 共题及高低年级共题，不是 503 道不同题。}

\key{主要判断：初赛的稳定骨架是分析、代数和解析几何；真正需要练的是把条件转换成估计、结构或几何关系。}数列、导数和积分经常交叉，矩阵题经常落到线性变换与多项式关系，几何题则需要主动选坐标、写参数、识别二次曲面。

\par\noindent\begin{tabularx}{\linewidth}{P{2.4cm}YY}
\toprule
优先级 & 应重点掌握 & 教材入口\\
\midrule
第一层：反复练熟 & 数列迭代与渐近；中值定理与导数估计；积分不等式与积分极限 & 数分 2.1、3.2--3.3、5.1；第 7 章的 1-B、2-G、3-A\\
第一层：反复练熟 & 矩阵秩、行列式、特征值、最小多项式、不变子空间与循环向量 & 高代几何第 1--3 章交叉阅读，特别是 1.3、2.3\\
第一层：单独准备 & 直线平面、距离与正交；切锥、轨迹、直母线、二次曲面 & 高代几何第 4 章基本工具、第 5 章、第 6 章\\
第二层：保证覆盖 & 级数敛散、幂级数；反常与含参积分；二次型、正定性 & 数分第 4、6 章；高代几何第 2 章\\
第三层：按目标加深 & Fourier 系统理论、特殊函数深入公式、病态构造、标准形算法的全部细节 & 本文逐节标出，不能整章一概略去\\
\bottomrule
\end{tabularx}

\para{读书策略}
先做一个代表题，再回教材补它所需的定理和两三个变式；同一方法的长串推广先留到第二轮。第 7 章不宜整体排到最后，它解释了许多初赛题的核心估计。解析几何也不宜只在考前背方程。

\para{适用范围}
默认以数学类初赛为主要备考场景，同时注明低年级决赛的加深方向。\key{高年级决赛的实分析、复分析、抽象代数、微分几何、概率和数值分析须另列准备。}这些内容在初赛样本中少，不代表在高年级决赛中不重要。

\clearpage
\topic{1\quad 材料范围、统计口径与可信边界}
\par\noindent\begin{tabularx}{\linewidth}{YrrY}
\toprule
范围 & 卷数 & 题位 & 用途\\
\midrule
常规初赛，第 1--17 届 & 24 & 150 & 阅读优先级的主要依据；第 11 届起有 A/B 卷\\
初赛补赛 & 6 & 36 & 检查方法是否稳健，单列避免第 13、14 届权重过大\\
第 1--4 届未分组决赛 & 4 & 30 & 补充早期综合题特征\\
第 5--17 届低年级决赛 & 13 & 119 & 观察证明深度、ODE 与综合题\\
第 5--17 届高年级决赛 & 13 & 168 & 识别高年级专业课模块，不混入初赛频率\\
\midrule
合计 & 60 & 503 & 含重复题位\\
\bottomrule
\end{tabularx}

\para{主要观察}
常规初赛 150 个题位分为：分析 77、代数 48、解析几何 25。第 11--17 届常规 A/B 卷以“3 个分析、2 个代数、1 个几何题位”为主体，第 11 届 B 卷另多 1 个分析题。\key{24 卷均有解析几何题}；第 17 届用四面体与向量体积，说明几何不能只准备曲面方程。

\para{统计口径}
\begin{itemize}
\item 每道大题记一个题位，小问不拆；交叉题只给一个主标签。题位数不等于分值，也不代表某个工具的全部使用次数。
\item 同届 A/B 卷按主题并集计算“覆盖届数”；共题不当作独立命题，补赛与决赛另列。阅读优先级还考虑方法的迁移性和前置作用，不直接按次数排序。
\item 历史频次用于决定阅读深浅，不视为下一届的题量或命题概率。
\end{itemize}

\para{材料与定位}
真题为你提供的个人解析汇编，题目归属沿用其标注；已核对代表页，未逐题审核全部答案。教程对照数学分析原书（409 页）及各章重排稿、高代几何全书重排稿（298 页）。\textbf{“真题 PDF 页”均从真题封面起数}；教材页码在各表另行说明。完整的 60 卷、503 题位分类仍保存在 \texttt{源码/资料提取/人工分类题目索引.json}，精简版不再逐卷列印。

\clearpage
\topic{2\quad 分析题型一：数列、迭代、极限与渐近}
\para{主要考点}
单调有界、Cauchy 准则、上下极限、Stolz 定理、Taylor 展开、Riemann 和；不动点、压缩估计、收敛速度及扰动递推。样本中对应主标签有 18 个常规初赛题位，覆盖 13 届。

\para{题目特征}
常见要求是“先证明极限存在，再求尺度或速度”，或者“给一个看似压缩的迭代，判断是否真能收敛”。仅由递推式解出候选不动点，再把它写成极限，通常不够。带参数时，临界值和边界点尤其需要证明。

\para{优先练会的三种切入}
\begin{enumerate}
\item \textbf{先找不变区间。}证明迭代始终留在一个区间，随后考虑单调性、奇偶子列或统一的压缩常数。导数绝对值处处小于 1，不自动给出全域统一的 $q<1$；常要先证轨道有界。
\item \textbf{在不动点处展开。}若 $x_n\downarrow0$，且
\[
x_{n+1}=x_n-cx_n^{1+\alpha}+o(x_n^{1+\alpha}),\qquad c,\alpha>0,
\]
则先验证条件，再用 $x_{n+1}^{-\alpha}-x_n^{-\alpha}\to\alpha c$，得到
$n x_n^\alpha\to1/(\alpha c)$。这里可迁移的是选取合适的倒数幂，不是记某道题的常数。
\item \textbf{把求和转成平均或积分。}卷积型求和先处理少量端点项，再在中间区间使用一致的小误差；不规则取样的 Riemann 和先对阶梯函数证明，再用上下和逼近。
\end{enumerate}

\exref{第 5 届初赛第 3 题，PDF 31：$x_{n+1}=f(x_n)$ 与 $nx_n$；第 16 届 B 初赛第 6 题，PDF 198：隐式递推与 $n x_n^\alpha$；第 14 届 A/B 初赛第 2 题，PDF 152/155：卷积极限。}
\exref{第 16 届 A 初赛第 6 题，PDF 195：两类函数的任意次复合；第 17 届低年级决赛解答第 5 题，PDF 217：有扰动、非全域统一压缩的递推。}

\para{教材怎样读}
数分 2.1 中“单调收敛、Stolz、定积分、压缩映射、上下极限、Taylor”这几组是主线，不宜只刷计算极限。第 7 章话题 3-A（微分与不动点）应紧接着读；先掌握例 7.3.4、7.3.6 的“轨道有界再收缩”思路。第 7 章 1-B 的差分估计则适合第二轮连接到渐近速度。

\clearpage
\topic{3\quad 分析题型二：导数、中值、凸性与极值}
\para{主要考点}
Rolle、Lagrange、Cauchy 中值定理；辅助函数与插值多项式；导数不等式、积分因子、Taylor 余项；凸凹性、驻点和边界极值。该主标签有 22 个常规初赛题位，覆盖 13 届，但其中不同子类型不能等量看待。

\para{三种反复出现的结构}
\begin{enumerate}
\item \textbf{多点条件换成零点条件。}题目给 $f,f'$ 的若干值，结论要求某个高阶导数取指定形式，优先构造 $F=f-P$，让 $P$ 同时吸收插值条件和目标导数，再反复用 Rolle。避免无目的地连续求导。
\item \textbf{全局存在性或增长限制。}给出 $f'$ 与 $f$ 的比较，结论说函数不存在、恒为零或不能超过某界，往往要选择 $e^{-cx}f$、$f^{1-\alpha}$、$f'/f$ 等量，把微分不等式积分化。
\item \textbf{凸性和极值。}一元题用导数单调与切线/弦线；多元题用 Hessian、紧集最值和边界分解。带区域约束的最小值不能只找内部驻点。
\end{enumerate}

\exref{第 1 届初赛第 7 题，PDF 7：三点共线与二阶导数零点；第 14 届补赛一 B 第 5 题，PDF 161：插值后反复 Rolle。}
\exref{第 8 届初赛第 5 题，PDF 65：$f'\ge f^\alpha$ 的全局不存在性；第 10 届初赛第 4、6 题，PDF 88--89：无限阶消失与导数控制；第 15 届 A/B 初赛第 2 题，PDF 177/180：三元函数驻点和约束极值。}

\para{一个最值得带走的估计}
若正函数在整个 $[0,\infty)$ 上满足 $f'\ge f^\alpha$，$\alpha>1$，则
\[
\bigl(f^{1-\alpha}\bigr)'=(1-\alpha)f^{-\alpha}f'\le-(\alpha-1).
\]
左侧原函数始终为正，却被迫在有限距离后变负，矛盾。这种“构造单调量并与全局定义域矛盾”的方法比记解方程更有用。

\para{教材怎样读}
数分 3.2、3.3 精读。3.3 的例 3.3.3、3.3.4、3.3.10 适合训练辅助多项式，例 3.3.8 适合识别全导数。第 7 章 1-B 的例 7.1.8--7.1.12、7.1.18，以及 2-G 的例 7.2.42--7.2.49 是导数估计与爆破/屏障的重点。

\para{不宜低估的基础}
Taylor 的余项阶、可导不推出导数连续、导数的 Darboux 性质，以及紧区间上最值存在。这些常决定证明是否成立。L'Hospital 是基础工具，但“把 3.4 全部技巧刷完”的收益通常低于上述几条主线。

\clearpage
\topic{4\quad 分析题型三：积分估计、矩与积分极限}
\para{主要考点}
Riemann 和、阶梯函数逼近、分部积分、变量代换、Cauchy--Schwarz/Hölder、凸性、等号与最优常数；积分中质量集中位置的比较。主标签有 16 个常规初赛题位，覆盖 11 届；许多导数题也以积分形式出现。

\para{四种常见题面与切入}
\par\noindent\begin{tabularx}{\linewidth}{P{3.3cm}Y}
\toprule
看到的条件 & 优先想到\\
\midrule
给 $\int f,\int xf$ 等矩，求 $\int f^2$ 的界 & 配方 $\int(f-P)^2\ge0$，把 $P$ 的系数按已知矩选好；或视作投影\\
给边界值、求导数积分控制函数积分 & 分部积分制造权函数，再用 Cauchy--Schwarz；等号由两函数成比例反推\\
$f^n$、$x^n$ 或归一化积分比值 & 找最大值点与局部衰减阶，分“主要邻域＋其余区域”，比较尺度\\
不规则选点、可积函数、振荡因子 & 先阶梯函数或简单函数，再用可积性控制误差；避免假设原函数处处连续\\
\bottomrule
\end{tabularx}

\exref{第 3 届初赛第 2 题，PDF 18：非负函数乘积与积分；第 5 届初赛第 5 题，PDF 32：偶函数、单调与凸性；第 12 届 B 初赛第 6 题，PDF 120：积分下界到平方积分。}
\exref{第 2 届初赛第 4 题，PDF 13：$n^2\int_0^1x^nf(x)\dd x$；第 9 届初赛第 6 题，PDF 78：两个端点同为最大值，但衰减阶不同；第 17 届 A 初赛第 6 题，PDF 210：具有自然密度的子序列 Riemann 和。}
\exref{第 17 届 A/B 初赛第 2 题，PDF 208/211：有限复指数和的二阶、四阶矩与 Hölder。它说明正交性值得掌握，但不能由一题推成 Fourier 全部理论都高频。}

\para{特别值得练的区别}
第 9 届题中 $f(x)=1-x^2+x^3$ 在 0、1 两端均为 1，但 $1-f(x)$ 在两端分别为 $x^2$ 阶和 $1-x$ 阶。积分质量分别约为 $n^{-1/2}$、$n^{-1}$，因此不是把两端的函数值简单平均。复习目标是\textbf{比较局部尺度}，不必升级成一整套精细渐近理论。

\para{教材怎样读}
数分 5.1 是重点章节。优先：例 5.1.13、5.1.15、5.1.19、5.1.24、5.1.31、5.1.35、5.1.39；真题 5.4.3--5.4.10。重排稿中保留了这些入口。第 7 章 1-A 与 2-A 分别补“微分/积分互换视角”和“极限能否传入积分”。

\clearpage
\topic{5\quad 分析题型四：级数、函数列与一致收敛}
\para{主要考点}
比较、积分判别、Dirichlet/Abel、分部求和、条件与绝对收敛、临界对数项；幂级数、生成函数、逐项运算及一致收敛。该主标签有 10 个常规初赛题位，覆盖 9 届。

\para{题目真正难在哪里}
常见级数不是一个标准通项直接代判别式，而是带有积分定义的系数、递推系数、取整造成的分块，或分母中有很小的扰动。交错符号本身不保证满足 Leibniz 的单调条件；“比值极限为 1”也不意味着判别已经结束。

\begin{enumerate}
\item \textbf{先做主项与余项分离。}例如 $\frac1{u_n+v_n}=\frac1{u_n}-\frac{v_n}{u_n^2}+\cdots$，需要检查余项级数是否可控。第 16 届 A 卷两题看似相近，但主尺度不同，不能只按交错形式判断。
\item \textbf{分块与分部求和。}对 $\lfloor\sqrt n\rfloor$ 或三角振荡，先按符号稳定的块处理；分部求和把振荡的部分和有界性与系数趋零连接起来。
\item \textbf{生成函数与递推。}从有理型幂级数读出递推，先证明系数关系，再求级数和。避免只计算前几项猜规律。
\item \textbf{一致收敛按最坏点控制。}先问区间是否无界、端点是否出问题；逐点估计要能对所有 $x$ 使用同一个 $N$，才能交换极限与积分等运算。
\end{enumerate}

\exref{第 9 届初赛第 2 题，PDF 76：单调数列定义的对数临界级数；第 10 届初赛第 5 题，PDF 88：含 $1/n$、$1/(n\ln n)$ 的比值条件；第 14 届 B 初赛第 5 题，PDF 156：有理函数系数与级数和。}
\exref{第 16 届 A 初赛第 5 题，PDF 195：取整分块与两个不同的分母尺度；第 1 届初赛第 4 题，PDF 6：导数统一有界、逐点收敛推一致收敛，且极限未必可导。}

\para{教材怎样读}
4.1 基本判别法、4.2 中分块/分部求和思想、4.3 中临界展开、4.4 积分估计、4.5.1--4.5.2 幂级数与生成函数、4.8 真题是主线。\textbf{Raabe、Bertrand/Gauss 至少知道条件和适用边界}，但 4.2 中层层高阶推广不用逐条背熟。

4.5.3 的 Abel/Cesàro 求和思想可选读；4.5.4 解析函数的完整延伸后置。原书 4.6 应先补三角正交与基本 Fourier 系数，而非从全部收敛判别定理开始。Bernstein 多项式详见后面的低优先级说明。

\clearpage
\topic{6\quad 分析题型五：反常、含参与多元分析}
\para{反常与含参积分的主要考点}
无穷远和瑕点分别比较；绝对/条件收敛；统一尾估计；参数求导、换序、对称代换。对应主标签有 9 个常规初赛题位，覆盖 8 届。频次低于导数和数列，但方法集中，不能放弃。

\para{常用结构}
含 $\log(1+s x^\beta)$、$\arctan(sx)$ 等参数时，先确认原积分有意义，再对参数求导化为有理积分，最后利用一个可计算的参数值定常数。参数靠近边界时，不能只说“显然可交换”。
\[
F(s)=\int_0^\infty K(x,s)\dd x,\qquad
F'(s)=\int_0^\infty \partial_sK(x,s)\dd x
\]
的关键是为参数紧区间找统一可积控制，或使用相应的一致收敛定理。对振荡尾部，Dirichlet 的有界原函数估计通常比绝对值粗估有效。

\exref{第 11 届 B 初赛第 2 题，PDF 103：反常积分的倒数代换与对称性；第 13 届 A 初赛第 5 题，PDF 133：振荡积分和换序；第 16 届 A/B 初赛第 2 题，PDF 193/196：对数含参积分。}

\para{教材怎样读}
数分 6.1、6.2 的敛散判断和 6.3 含参积分优先；例 6.3.1 的参数边界、6.3.4--6.3.6 的一致收敛、6.3.7--6.3.11 的计算路径值得练。6.4 先掌握 $\Gamma$、$B$ 定义、递推、$\Gamma(1/2)$ 及基本换元；复杂反射、倍增公式和对数 Gamma 的一长串应用可后置。

\para{多元微分与曲线曲面积分：出现较少但须保留基础}
初赛样本中独立的常规曲线/曲面积分计算题少于一元估计题；多元极值和 Laplacian 极值论证仍会出现。决赛填空题则经常出现重积分、曲线与曲面积分，不能把“初赛较少”带入决赛复习。
\begin{itemize}
\item 多元极值：Hessian、约束和边界；出现 $\Delta u$、边界值时，先用内部极大点处梯度为零、Hessian 半负定的基本事实。
\item 曲线/曲面积分：参数化、方向、对称性、补面或挖去奇点；Green/Gauss/Stokes 的条件应会核对。
\end{itemize}
\exref{第 1 届初赛第 6 题，PDF 7：圆盘积分与 Green；第 13 届 A/B 初赛第 2 题，PDF 132/135：$\Delta(u^2+v^2)$ 与极大值；第 15 届 A/B 初赛第 2 题，PDF 177/180：边界极值。}
\para{对应阅读}
数分第 3 章基本内容、5.2--5.3 的基本公式和少量代表题必须覆盖；纯粹大量参数积分计算可第二轮再刷。这里的“少”是相对初赛的独立题位，不是排除某个必需工具。

\clearpage
\topic{7\quad 代数主线：从矩阵运算进入线性结构}
\para{主要考点}
特征值与特征空间、相似与 Jordan、最小多项式、不变子空间、核像分解、循环向量、交换矩阵、矩阵多项式和矩阵方程。常规初赛中，“线性空间/循环结构/矩阵方程”主标签有 15 题位、覆盖 12 届；“谱/相似/标准形”另有 8 题位、覆盖 7 届。\textbf{两类经常共用谱理论，不能把标签次数当作知识点的全部出现次数。}

\para{反复出现的核心对象}
第 1 届初赛伴随型矩阵的交换子空间、第 5 届低年级决赛 Hessenberg 型矩阵的交换子空间、第 15 届初赛的循环向量、第 17 届 A 初赛的向量最小多项式与交换矩阵，都属于同一条结构主线：\key{一个向量的迭代轨道能否生成整个空间。}

若 $\alpha,A\alpha,\ldots,A^{n-1}\alpha$ 是一组基，且 $AB=BA$，写
$B\alpha=f(A)\alpha$，则
\[
BA^k\alpha=A^kB\alpha=f(A)A^k\alpha.
\]
于是 $B=f(A)$。此证明解释了为什么交换矩阵由一个列向量决定，也把维数、最小多项式、循环基联系起来。这比只背“可交换所以可同时对角化”可靠；后一句缺少条件时是错误的。

\exref{第 1 届初赛第 2、3 题，PDF 5--6；第 5 届低年级决赛第 3 题，PDF 35；第 15 届 A/B 初赛第 3 题，PDF 178/181；第 17 届 A 初赛第 4 题，PDF 209。}

\para{做题时优先问的五个问题}
\begin{enumerate}
\item 条件给了哪条多项式关系 $p(A)=0$？能否把高次幂降到低次？
\item 哪些子空间自然不变：$\ker p(A)$、$\operatorname{Im}p(A)$、某个向量的循环子空间？
\item 换一组适应这些子空间的基后，是否得到分块上三角或分块对角？
\item 所求是相似不变量吗：秩、迹、行列式、特征多项式、最小多项式？
\item 数域是否影响结论：实数、复数、有理数上是否能分解或取到特征向量？
\end{enumerate}

\para{教材怎样读}
高代几何第 1 章看运算、不变量与分块；第 2 章要读线性变换和空间分解，不只读二次型；第 3 章补多项式工具。具体优先：例 1.3.17、1.3.20，例 2.3.2、2.3.15--2.3.17。例 2.3.17 在重排 PDF 136--137 页，直接对应第 15 届循环向量题，是很好的精读点。

\clearpage
\topic{8\quad 代数补全：秩、行列式、二次型与多项式}
\para{秩、行列式与构造：不能只练展开}
主标签有 11 个常规初赛题位，覆盖 6 届。常用工具包括秩分解、核像、分块行列式、秩一更新、伴随矩阵按秩分情形，以及适当构造矩阵实现给定条件。
\begin{itemize}
\item 遇 $uv^\mathsf T$，先识别秩一结构；$I+uv^\mathsf T$ 的可逆性由 $1+v^\mathsf Tu$ 决定。
\item 遇高维、特殊取值和存在性，先降到一个小块或从多线性出发；不要被 2024、2029 等年份数字引向大规模计算。
\item 遇伴随矩阵，先分 $\rank A\le n-2$、$\rank A=n-1$、可逆三种结构，再使用 $AA^*=\det(A)I$。
\end{itemize}
\exref{第 15 届 B 初赛第 4 题，PDF 181：选择对角 $\pm1$ 使矩阵可逆；第 16 届 A/B 初赛第 3 题，PDF 194/197：伴随矩阵；第 16 届 B 第 4 题，PDF 197：秩一更新；第 17 届 A/B 第 3 题，PDF 209/212：给定行列式的构造。}

\para{二次型、正定与正交：基础必须完整}
主标签有 7 个常规初赛题位，覆盖 6 届；在决赛填空与证明中还反复出现。重点是实对称矩阵的谱定理、正交对角化、合同与惯性、正定判别，以及把 $x^\mathsf TAx$ 改写成平方和/积分平方。非对称实矩阵的二次型只取决于 $(A+A^\mathsf T)/2$。

必须分清：相似变换与合同变换；标准形与规范形；代数重数与几何重数。对非对称矩阵不能直接套实对称矩阵的谱结论。
\exref{第 4 届初赛第 7 题，PDF 27：方程有解、相似与合同的区别；第 10 届初赛第 2 题，PDF 87：由元素控制判定二次型；第 11 届 A 第 4 题，PDF 101：二次型组成的空间。}

\para{多项式：有稳定用途，也有较尖的独立题}
主标签有 7 个常规初赛题位，覆盖 5 届。Bezout、互素、根与重数、Vieta、最小多项式、Cayley--Hamilton 是第一层；整系数根的约束、专门的不可约证明和更深的数论构造为第二层。
\[
\gcd(p,g)=1 \ \Longrightarrow\ up+vg=1
\ \Longrightarrow\ v(A)g(A)=I\quad\text{当 }p(A)=0.
\]
这是“多项式互素”和“矩阵可逆”之间应当熟练使用的桥梁。
\exref{第 12 届 B 初赛第 3 题，PDF 120；第 13 届 A/B 第 3 题，PDF 132/135：整系数多项式不能全实根；第 14 届 A 第 4 题，PDF 153：不可约性。}
\para{对应阅读}
高代几何 1.2、1.3，2.2、2.3 和 3.2、3.3。第 3 章不应整体降级；优先矩阵多项式与整除工具，再读专门技巧。Smith/有理标准形全套算法可后置，Jordan 与最小多项式的结构作用不可后置。

\clearpage
\topic{9\quad 解析几何：稳定的一题，不能只背曲面名称}
\para{主要考点}
点积、叉积、混合积和体积；直线平面、异面公垂线；球面、柱面、锥面、旋转面；二次曲面分类、截线、直母线、切线与切点平面、动点轨迹。\textbf{常规初赛 24 卷均有几何题，覆盖全部 17 届。}

\par\noindent\begin{tabularx}{\linewidth}{P{2.8cm}YY}
\toprule
题型 & 常用切入 & 代表定位\\
\midrule
向量、体积、重心 & 用混合积表达有向体积；以一顶点为原点，建立适应四面体的基 & 第 17 届 A/B 第 1 题，PDF 208/211\\
异面直线与公垂线 & 设两线上的参数点，令连线与两个方向都正交，解二元方程 & 第 6、7 届初赛第 1 题，PDF 42/53\\
轨迹与曲面识别 & 把切球、距离等几何条件转换成距离等式；充分性也要验证 & 第 11 届 A 第 1 题，PDF 100；第 15 届 A 第 1 题，PDF 177\\
切锥与切点平面 & 直线代入二次式，相切令判别式为零；再提取极平面 & 第 13 届 B 第 1 题，PDF 135；第 16 届 A 第 1 题，PDF 193\\
直母线与旋转面 & 直线参数代入曲面，比较各次系数；旋转消元要检查退化情形 & 第 10、13、14、16 届相关初赛首题\\
\bottomrule
\end{tabularx}

\para{最通用的二次曲面工具}
把曲面写成 $Q(X)=X^\mathsf THX+2b^\mathsf TX+c=0$。直线 $X=P+tv$ 代入后为
\[
(v^\mathsf THv)t^2+2v^\mathsf T(HP+b)t+Q(P)=0.
\]
相切研究重根；整条直线落在曲面上要求三个系数都为零；截线形状研究限制到平面后的二次型。选对坐标后，许多空间几何题都能用这一框架解决。

\para{教材怎样读}
高代几何第 4 章的工具先过一遍，第 5 章选做方程、距离、公垂线，第 6 章重点读。精读入口：例 5.3.2--5.3.3；例 6.2.4、6.2.10、6.2.19--6.2.20；例 6.3.6--6.3.8。这些覆盖了公垂线、动球轨迹、直母线与截面。

\para{相对少考的内容怎样处理}
心脏线反演等特殊曲线是样本中的少数题，保留极坐标、参数化、反演定义的基本认识即可，先不收集大量冷门曲线。第 4 章独立的长向量恒等式证明题较少，但其运算工具和重心体积不能跳过，第 17 届正好说明这一点。

\clearpage
\topic{10\quad 数学分析教程：按节安排阅读重点（上）}
以下节号沿原书。原书正文页 1 对应原 PDF 第 12 页；表中给出\textbf{原书印刷页码}，避免与各章重排页码混淆。重排版没有完整覆盖原书，缺项见第 13 节。

\begingroup\small\renewcommand{\arraystretch}{1.13}
\begin{longtable}{P{3.0cm}P{1.5cm}P{10.9cm}}
\toprule
章节与原书页 & 优先级 & 阅读重点及可后置内容\\
\midrule\endhead
第 1 章 Euclid 空间，1--10 & 基础必会；短读 &
紧致、子列、连续函数最值、一致连续、介值和集合直径必须能用。开集分解等完整拓展证明可按基础情况选读。样本独立考集合性质的题较少，但紧致是后续估计的重要前置。\\
2.1 极限，11--37 & 精读 &
按单调、Stolz、积分、级数、压缩、上下极限、Taylor 分组练习。每组先掌握一个能迁移的证明，再做变式。少花时间重复同一套机械等价替换。\\
2.2 连续性，38--57 & 精读核心 &
连续与一致连续、紧致性、Lipschitz、反函数、不动点、可导/连续之间的反例。复杂病态函数的完整构造可后置，但不能把概念边界一并删掉。\\
2.3 典型性质，58--61 & 第一轮后置 &
Baire 类、“典型”连续函数、无处单调/可导的主剩集证明。常规初赛未见直接要求整套理论的题；不因此降低普通连续性和构造反例的优先级。\\
2.4 真题，62--71 & 精读 &
做完后写出每题的关键估计和所需条件。书中旧题与本汇编重合，不算另一份独立的频率证据。\\
3.1--3.3，72--93 & 精读 &
中值定理、插值辅助函数、凸性、导数与函数增长。多元微分和 Taylor 基础保留；优先本文第 3 节列举的代表题。\\
3.4 L'Hospital，94--96 & 会用即可 &
熟悉适用条件，能处理必要的极限；细碎技巧不宜挤占中值和估计训练。重排主文件未包含本节，须确认自己基础已覆盖。\\
3.5 真题，97--115 & 精读 &
与第 7 章导数估计联读；做题时明确局部条件与整个定义域条件各在哪里用。重排文件标题为“节选”，不假定保留了所有原题。\\
4.1--4.4，116--144 & 精读主线，推广选读 &
常用判别、分块、分部求和、临界比值展开、积分估计必须掌握。高阶 Raabe/Kummer 等成串推广只选典型，不要求全部记忆。\\
4.5 幂级数，145--149 & 基础与生成函数精读 &
收敛区间、逐项微积分、从有理函数建立递推。Abel/Cesàro 求和选读；解析函数的完整理论放在复分析目标下准备。\\
\bottomrule
\end{longtable}\endgroup

\clearpage
\topic{11\quad 数学分析教程：按节安排阅读重点（下）}
\begingroup\small\renewcommand{\arraystretch}{1.13}
\begin{longtable}{P{3.0cm}P{1.5cm}P{10.9cm}}
\toprule
章节与原书页 & 优先级 & 阅读重点及可后置内容\\
\midrule\endhead
4.6 Fourier，150--165 & 基本工具补读；系统理论选读 &
优先三角/复指数正交、基本系数、Parseval/Bessel 的含义和适用条件。第 17 届初赛用到正交与矩估计，第 8 届决赛出现 Fourier 变换；不能把全节说成“没用”。Dini-Lipschitz、Dirichlet-Jordan 的全部证明、Schwartz 变换全套第一轮可后置。\\
4.7 Bernstein，166--172 & 第一轮后置 &
先知道以多项式逼近连续函数的思想，及概率/方差控制误差的基本路径。本次常规初赛未见必须直接调用 Bernstein 构造的独立题，不必优先练高阶导数逼近与一串推广。\\
4.8 真题，173--190 & 精读 &
以敛散判断、函数列与逐项运算为主。碰到更深的 Fourier/逼近题再回补原书对应理论。\\
5.1 Riemann 积分，191--257 & 最高优先组之一 &
不等式、矩、分部积分、积分极限、Riemann 和、阶梯逼近、积分方程。代表题见本文第 4 节；复杂重复计算与高维专门公式第二轮选做。\\
5.2--5.3，258--283 & 基础必会；计算选做 &
两型曲线/曲面积分、参数化、方向、Green/Gauss/Stokes、补面/挖奇点。初赛独立计算较少，决赛填空要提高熟练度。\\
5.4 真题，284--297 & 精读 &
把每题归到估计、换序、局部集中、方程或积分公式，避免只记答案。例 5.4.9 的双端点尺度比较尤其值得读。\\
6.1--6.3，298--323 & 精读 &
无穷远与瑕点分别判断、振荡积分的尾估计、含参连续性和求导换序。先学如何证明“可交换”，再熟悉经典计算。\\
6.4 Euler 积分，324--330 & 基本公式必会，深入后置 &
$\Gamma$、$B$ 与换元、递推是通用工具；对数 Gamma、反射/倍增公式的长推演可按题补读。\\
6.5 真题，331--340 & 精读 &
估计与微分不等式、反函数几何、振荡参数、函数列迭代。不是只练反常积分求值。\\
7.1，341--348 & 1-A、1-B 精读 &
积分/微分对应、导数与差分估计，先读例 7.1.8--7.1.12、7.1.18、7.1.25。1-C 中矩阵序、有限维函数空间的推广选读。\\
7.2，349--375 & 2-A、2-G 精读 &
极限与积分换序、积分因子、屏障与爆破。2-C 中例 7.2.17--7.2.20 的积分平方/正定矩阵补读；2-D 的全套广义中值拓展可后置。\\
7.3，376--394 & 3-A 精读，3-B/3-C 选读 &
不动点、轨道有界、一致收敛、隐函数。3-D 的题目再加工与阶乘三角级数专题先放第二轮，掌握方法后再扩展。\\
\bottomrule
\end{longtable}\endgroup

\clearpage
\topic{12\quad 高等代数与解析几何教程：怎样精读}
下表页码均为\textbf{《高等代数与解析几何\_全书重排稿.pdf》的物理页码}。该书共 298 页，正文六章；节号和例题号仍可沿用原书。

\begingroup\small\renewcommand{\arraystretch}{1.15}\par\noindent\begin{tabularx}{\linewidth}{P{3.3cm}Y}
\toprule
章节、页码与优先级 & 重点与取舍\\
\midrule
第 1 章，PDF 7--92\newline 矩阵与线性变换\newline 精读 &
矩阵运算、分块、秩一结构、秩/迹/行列式、特征值、相似与 Jordan、矩阵方程。1.2 选典型方法，1.3 真题优先。例 1.2.3（秩一矩阵幂，PDF 9--10）、1.3.17（交换矩阵，PDF 61--62）、1.3.20（特征向量刻画数量算子，PDF 65）值得精读。复杂特殊行列式的重复算例选做。\\
第 2 章，PDF 93--165\newline 线性方程组与二次型\newline 精读 &
核像与维数、空间分解、正交/酉变换、二次型、正定性、线性变换的结构。例 2.2.11（PDF 103）、2.2.17（PDF 108），及 2.3.2（PDF 117 起）、2.3.12（PDF 132）、2.3.15（PDF 134 起）、2.3.17（PDF 136 起）是有效入口。不要按章名误以为这里只是解线性方程组。\\
第 3 章，PDF 166--200\newline 多项式与矩阵方程\newline 工具精读，尖题选读 &
互素与 Bezout、零化多项式、矩阵可逆、根与重数、矩阵方程有无解。先读 3.2 和 3.3，再按需要补拓展训练。一般不可约技巧和复杂整系数构造排在这些通用工具之后。\\
第 4 章，PDF 201--217\newline 向量代数\newline 工具必会，长恒等式选读 &
点积、叉积、混合积、Lagrange 恒等式、共面、距离和体积。第 17 届首题提示：这些工具能独立成题。若已熟练，快速复习即可，不需每个恒等式都反复证明。\\
第 5 章，PDF 218--245\newline 平面与直线\newline 核心计算练熟 &
法向与方向、相交/平行/垂直、异面距离、公垂线、参数消元。5.2 选基础代表题，5.3 全部三道真题是很短而有效的一组；例 5.3.2、5.3.3 位于 PDF 237--238。\\
第 6 章，PDF 246--297\newline 二次曲线与二次曲面\newline 解析几何重心 &
球/柱/锥/旋转面、切线与极平面、直母线、轨迹、平面截线。例 6.2.4（PDF 250）、6.2.10（PDF 255 起）、6.2.19（PDF 265）、6.2.20（PDF 266 起）、6.3.6--6.3.8（PDF 287--289）优先。特殊曲线和大量同构计算可选做。\\
\bottomrule
\end{tabularx}\endgroup

\para{书内栏目顺序}
每章先快速过“竞赛要点与难点”，再做“真题选讲”，不会的回看“范例”。“能力拓展与训练”用于补某种方法，暂不逐题做完；“训练全解”在独立尝试后用来核对关键步骤，不按答案顺序阅读。

\para{版本说明}
高代几何重排稿用局部原图保留许多复杂公式，文字提取可能丢公式，部分识别文字也不完美；遇到不清楚的式子应看版面或配套原书。本报告的例题号按重排版实际保留的编号定位，不把局部原图缺乏可提取文本当作书中缺内容。

\clearpage
\topic{13\quad “考得少”与“重排缺项”要分开判断}
\par\noindent\begin{tabularx}{\linewidth}{P{3.2cm}P{3.3cm}Y}
\toprule
内容 & 本材料中的判断 & 复习处理\\
\midrule
2.3 连续函数典型性质、Baire 主剩集 & 常规初赛中未见直接要求整套理论 & 后置完整证明；保留普通紧致与连续性基础\\
4.2 高阶判别法、4.3 特殊数字串与等分布拓展 & 有方法价值，具体高级版本直接题位少 & 掌握比较、分部求和及临界项，推广选一两题\\
Fourier 系统收敛理论、Bernstein 高阶逼近 & 初赛直接系统考查少；正交工具有明确用途 & 补正交、基本展开和误差思想；后置整套技术证明\\
Gamma 函数深入恒等式与冷门积分 & 本次初赛中不是主要独立主线 & 基本 $B/\Gamma$ 保留，其他随题补\\
曲线曲面积分的繁复计算 & 相对初赛一元估计题少，决赛填空仍常见 & 基本公式与方向必会，初赛先做代表题\\
Smith/有理标准形完整算法 & 与循环、Jordan、最小多项式相比，独立机械计算较少 & 结构思想优先；不由此跳过 Jordan 与多项式\\
长向量恒等式、特殊曲线反演 & 专门长题少，但基本向量工具不可少 & 保留距离、体积与参数化；特殊构造后置\\
\bottomrule
\end{tabularx}

\para{现有数分重排版的实际缺项}
\begin{enumerate}
\item 本次未找到第 1 章独立重排文件；第 2 章主文件包含 2.1、2.2、2.4，没有 2.3。第 3 章未包含 3.4，3.5 标为“节选”；第 4 章未包含 4.6、4.7。
\item 第 5 章开头明确略去常规性质、繁复定义与重复计算，保留代表题及真题；这是精简稿，不能替代全部定义条件。第 6 章也应与原书按需要对照，而不按篇幅判断重要性。
\item 第 7 章明确省去 7.2 的 2-B、2-C、2-D、2-F，及例 7.3.11--7.3.13；证明做了精简，并注明不录入反例。
\end{enumerate}

\para{应主动补看的最小集合}
优先补\textbf{概念边界的反例}（逐点/一致、连续/可导、函数趋零/导数趋零）、\textbf{原书 4.6 的正交工具}，以及\textbf{原书 7.2 话题 2-C 的例 7.2.17--7.2.20}（印刷页 361--362，原 PDF 372--373：用积分表示证明矩阵正定）。这些被略去，不代表低价值。

\para{不能套用“初赛少考”的高年级内容}
第 5--17 届高年级决赛卷逐卷出现了实分析、复分析、抽象代数、微分几何、概率相关题位；数值方法也逐卷出现相关题位。这些模块在高年级范围中有稳定位置，两册教程并不完整覆盖，须以相应课程教材另补。本文不把它们列为所有 CMC 考生都可略过的内容。

\clearpage
\topic{14\quad 按目标使用这份报告}
\para{如果以初赛为主}
先走共同主线：数分 2.1、3.2--3.3、5.1、6.1--6.3，配合第 7 章 1-B、2-A、2-G、3-A；代数用第 1--3 章串起秩、谱、空间和多项式；几何用第 4 章工具加第 5--6 章。随后用第 4 章级数和二次型补齐覆盖。A/B 的共同题很多，不能把 B 类简单理解成“只算不证”。

\para{如果准备低年级决赛}
在上述基础上，提高填空计算的完整性，补重积分、线面积分与 ODE。样本中低年级决赛常出现一阶/二阶方程、全局存在与唯一性、周期解、非唯一拼接、微分不等式。\textbf{第 7 章只有相关入口，不能代替完整的 ODE 基础。}

\para{如果准备高年级决赛}
共同基础题继续练，同时分别准备六类专业课模块。实分析看测度、收敛定理和 $L^p$；复分析看 Cauchy、最大模、Schwarz、零点与留数；抽代看群环域的结构证明；微分几何看曲率、基本形式、测地线；概率看分布、条件期望、极限与随机过程；数值看迭代谱半径、逼近误差、ODE 格式与稳定性。这里是从样本题目归纳的方向，不是当前官方大纲的替代文本。

\para{第一轮代表题：用来暴露方法缺口}
\par\noindent\begin{tabularx}{\linewidth}{P{2.8cm}YY}
\toprule
主题 & 代表题（真题 PDF 页） & 做完应能回答\\
\midrule
迭代与速度 & 5 初 3（31）；16 B 初 6（198） & 存在性与速度各如何证明？\\
中值与辅助函数 & 1 初 7（7）；14 补一 B 5（161） & 目标导数如何决定辅助多项式？\\
导数与全局估计 & 8 初 5（65）；10 初 4（88） & 哪个量单调，定义域在哪里用？\\
积分极限与不等式 & 2 初 4（13）；9 初 6（78）；12 B 初 6（120） & 主质量在哪里，权函数如何选？\\
级数与参数积分 & 10 初 5（88）；16 A 初 5（195）；16 A/B 初 2（193/196） & 临界项是什么，何时可换序？\\
线性结构 & 1 初 2（5）；15 A/B 初 3（178/181）；17 A 初 4（209） & 循环向量与交换子空间如何联系？\\
秩与多项式 & 12 B 初 3（120）；16 B 初 4（197）；17 A/B 初 3（209/212） & 能否降到低维或互素关系？\\
解析几何 & 7 初 1（53）；15 A 初 1（177）；16 A 初 1（193）；17 A 初 1（208） & 坐标怎么选，必要条件是否充分？\\
\bottomrule
\end{tabularx}

\para{防止过拟合的做法}
同一方法做会后，换一种题面验证；A/B 共题和书中重复真题只算一次学习经历。保留第 16--17 届若干整卷作限时检查，不用它们反复调参；以“能否独立写出证明与条件”决定回读章节，而不是以押中某个年份的形式决定投入。

\end{document}
